% !TeX root = ../main.tex
% 1. fejezet — Bevezetés
% A ftsrg/thesis-template-latex sablonba \include-dal illesztendő tartalomfájl.
% Hivatkozott források: Backend/HUNting.API/Models/Users/AccountType.cs,
%   thesis/planning/feature_matrix.md (funkciókör), Frontend/HUNting/src/routes/appRouter.tsx
% MEGJEGYZÉS: az 1.4 alfejezet konkurens megoldásainak konkrét megnevezése a szerzővel
%   egyeztetendő (ld. TODO-kommentek) — jelenleg kategóriaszinten, óvatos állításokkal íródott.

\chapter{Bevezetés}
\label{chap:bevezetes}

\section{Motiváció és problémafelvetés}
\label{sec:bev-motivacio}

A vadászat Magyarországon szervezett keretek között, jellemzően vadásztársaságokban folyó
tevékenység, amelynek mindennapjait sokrétű szervezési és adminisztrációs feladat kíséri. A
társaságoknak kommunikálniuk kell tagjaikkal a közös vadászatokról, munkákról és hírekről;
kezelniük kell a tagfelvételt; a vadászok eszközöket -- fegyvert, optikát, ruházatot,
felszerelést -- adnak-vesznek egymás között; a területi tájékozódáshoz megfigyelőpontok, etetők
és területhatárok ismerete szükséges; emellett a vadász saját vadászati eseményeiről is
nyilvántartást kíván (vagy köteles) vezetni.

Ezek az igények ma jellemzően egymástól független, az adott célra csak részben alkalmas
csatornákon oszlanak szét: a társasági kommunikáció telefonon, levelezőlistákon vagy közösségi
médiás csoportokban zajlik; az adásvétel általános célú hirdetési oldalakon vagy zárt
csoportokban történik, ahol a vadászati eszközök sajátos kategóriái nem jelennek meg; a területi
tudás gyakran papírtérképeken vagy egyszerűen a tagok fejében él; a naplóvezetés pedig
papíralapú vagy táblázatkezelős. A széttagoltság többszörös adatrögzítést, információvesztést
és belépési küszöböt jelent -- különösen az új tagok számára, akik a közösség hallgatólagos
tudásához nehezen férnek hozzá.

A dolgozat tárgya egy olyan webalkalmazás -- a \emph{HUNting} vadászati közösségi és hirdetési
portál -- megtervezése és megvalósítása, amely ezeket az igényeket egyetlen, integrált
rendszerben szolgálja ki: közös felhasználókezeléssel, egymásra épülő funkciókkal és egységes,
magyar nyelvű felülettel.

\section{A vadászati terület digitalizációs lehetőségei}
\label{sec:bev-digitalizacio}

A vadásztársadalom digitális eszközhasználatát kettősség jellemzi. Egyfelől a terület
hagyományőrző: a közösségi élet személyes jelenlétre épül, az adminisztráció sok helyen ma is
papíralapú. Másfelől az általános célú digitális csatornák -- közösségi médiás csoportok,
üzenetküldő alkalmazások, apróhirdetési portálok -- használata a vadászok körében is
mindennapos, ami jelzi, hogy a digitális eszközökkel szembeni ellenállás nem elvi, hanem
kínálati kérdés: az általános célú eszközök nem képezik le a terület sajátosságait.

Egy szakterületre szabott portál hozzáadott értéke éppen ebben rejlik. A tartományspecifikus
fogalmak -- vadásztársaság, tagsági kérelem, társasági szerep, vadfaj, megfigyelőpont, etető --
első osztályú entitásként jelenhetnek meg benne, nem pedig szabad szöveges bejegyzésekként egy
általános platformon. Az integráció további előnye az adatok összekapcsolhatósága: ugyanaz a
fiók ad fel hirdetést, igényel tagságot és vezet naplót; a társasági tagság jogosultságot
adhat közös térképi rétegekhez; a hirdetés tulajdonosa lehet maga a társaság is. Ezek a
kapcsolatok a széttagolt megoldásokban nem, vagy csak kézi egyeztetéssel valósulnak meg.

\section{Az alkalmazás céljai és célcsoportjai}
\label{sec:bev-celok}

A HUNting két célcsoportot szolgál ki, amelyeket a rendszer két, jogosultságaiban és
felületeiben eltérő fióktípusként kezel (a backend \texttt{AccountType} felsorolt típusa
szerint): az \emph{egyéni vadászokat} (Person) és a \emph{vadásztársaságokat} (Society).
E kettősség a rendszer egészén végigvonul, a regisztrációtól a menüszerkezetig
(\ref{chap:felhasznalok}.~fejezet). A funkcionális célokat hat funkcióterület valósítja meg,
amelyek egyben a dolgozat gerincfejezeteit is kijelölik:

\begin{enumerate}[nosep]
	\item \emph{felhasználókezelés, jogosultságok és adminisztráció} -- regisztráció a két
	      fióktípussal, bejelentkezés, profil, admin-felület (\ref{chap:felhasznalok}.~fejezet);
	\item \emph{hirdetési piactér} -- vadászati eszközök adásvétele kategóriákkal, szűréssel,
	      képfeltöltéssel (\ref{chap:hirdetesek}.~fejezet);
	\item \emph{vadásztársaságok és tagságkezelés} -- társaságkatalógus, tagfelvételi
	      kérelmek életciklusa, tag- és bejegyzéskezelés (\ref{chap:tarsasagok}.~fejezet);
	\item \emph{vadásznapló} -- személyes vadászati események rögzítése képekkel
	      (\ref{chap:naplo}.~fejezet);
	\item \emph{térképes funkciók} -- megfigyelőpontok, etetők, veszélyhelyek és területek
	      személyes és társasági rétegeken (\ref{chap:terkep}.~fejezet);
	\item \emph{AI-chatbot és időjárás} -- vadászati asszisztens és időjárás-tájékoztatás
	      (\ref{chap:chatbot}.~fejezet).
\end{enumerate}

Nem funkcionális célként a rendszer magyar nyelvű, reszponzív (asztali és mobil nézetben
egyaránt használható -- \ref{fig:home-responsive}.~ábra a~\ref{chap:alapok}.~fejezetben),
valamint szerveroldalon kikényszerített jogosultsági modellre épül. Fontos a célok őszinte
lehatárolása is: a rendszer demonstrációs érettségű -- éles üzemeltetésre, terheléses
méretezésre és teljes biztonsági keményítésre nem készült fel; az ismert korlátokat minden
funkciófejezet, összegzésüket a \ref{chap:teszteles}.~és a \ref{chap:osszefoglalas}.~fejezet
tartalmazza.

\section{Hasonló megoldások}
\label{sec:bev-hasonlo}

% TODO: a konkrét konkurens megoldások (nemzetközi vadásznapló-appok, hazai hirdetési
% portálok) megnevezése és jellemzése a szerzővel egyeztetendő; az alábbi szöveg
% kategóriaszinten, termékmegnevezés nélkül írja le a mezőnyt. Konkrét termékekről csak
% publikus, ellenőrizhető információ alapján szabad állítást tenni.

A HUNting által lefedett igényekre léteznek részmegoldások, ezek azonban jellemzően a terület
egy-egy szeletét célozzák. A \emph{vadásznapló-alkalmazások} (elsősorban nemzetközi mobil
appok) a vadászati események rögzítésére és részben térképes megjelenítésére szolgálnak,
közösségi és hirdetési funkciók nélkül, és magyar lokalizációjuk többnyire hiányzik. Az
\emph{általános hirdetési portálok} a vadászati eszközök adásvételét is hordozzák, de a
tartományspecifikus kategorizálás és a vadászközösségi kontextus nélkül; a fegyverhirdetések
kezelése ráadásul speciális szabályozási környezetet igényel. A \emph{közösségi médiás
csoportok} a társasági kommunikáció leggyakoribb mai csatornái: belépési küszöbük alacsony,
cserébe strukturálatlanok -- tagfelvételi folyamat, szerepek, naplózás vagy térképi rétegek
nem fejezhetők ki bennük natívan.

A HUNting elhatárolása e mezőnytől az \emph{integráció}: nem egy-egy részfunkcióban kíván a
specializált megoldásoknál többet nyújtani, hanem a vadászati közösségi élet funkcióit kapcsolja
össze egyetlen, közös fiók- és jogosultsági modellre épülő, magyar nyelvű rendszerben. A
választott felépítés sajátossága a kettős fióktípus: a társaság nem csupán címke vagy csoport,
hanem saját fiókkal, felülettel és jogosultságokkal rendelkező szereplő.

\section{A dolgozat felépítése}
\label{sec:bev-felepites}

A dolgozat szerkezete a rendszer funkcióközpontú jellegét követi. A keretfejezetek közül
a~\ref{chap:technologiak}.~fejezet a felhasznált technológiákat és választásuk indokait,
a~\ref{chap:architektura}.~fejezet a háromrétegű architektúrát és a rétegek közötti
kommunikációt, a~\ref{chap:alapok}.~fejezet pedig a minden funkció által közösen használt
infrastruktúrát -- adatbázis-kontextust, backend-csővezetéket, kliensvázat -- mutatja be.

A dolgozat gerincét a \ref{chap:felhasznalok}--\ref{chap:chatbot}.~funkciófejezetek adják.
Mindegyikük azonos belső sablont követ: a funkció céljának és felhasználói működésének
bemutatása után az adatmodellt, a backend-, majd a frontend-megvalósítást tárgyalja, végül a
teljes adatfolyamot, a biztonsági megfontolásokat és -- őszinte felmérésként -- az ismert
korlátokat összegzi. Ez a felépítés teszi lehetővé, hogy minden funkció a teljes vertikumában,
az adatbázistól a felhasználói felületig egy helyen legyen olvasható.

A gerincet két lezáró fejezet követi: a~\ref{chap:teszteles}.~fejezet a háromszintű
tesztkészletet és a futtatási környezetet dokumentálja, a~\ref{chap:osszefoglalas}.~fejezet
pedig az eredményeket összegzi és a továbbfejlesztési irányokat jelöli ki.
